% Options for packages loaded elsewhere
% Options for packages loaded elsewhere
\PassOptionsToPackage{unicode}{hyperref}
\PassOptionsToPackage{hyphens}{url}
%
\documentclass[
  english,
  russian,
  12pt,
  a4paper,
  DIV=11,
  numbers=noendperiod]{scrreprt}
\usepackage{xcolor}
\usepackage{amsmath,amssymb}
\setcounter{secnumdepth}{5}
\usepackage{iftex}
\ifPDFTeX
  \usepackage[T1]{fontenc}
  \usepackage[utf8]{inputenc}
  \usepackage{textcomp} % provide euro and other symbols
\else % if luatex or xetex
  \usepackage{unicode-math} % this also loads fontspec
  \defaultfontfeatures{Scale=MatchLowercase}
  \defaultfontfeatures[\rmfamily]{Ligatures=TeX,Scale=1}
\fi
\usepackage{lmodern}
\ifPDFTeX\else
  % xetex/luatex font selection
\fi
% Use upquote if available, for straight quotes in verbatim environments
\IfFileExists{upquote.sty}{\usepackage{upquote}}{}
\IfFileExists{microtype.sty}{% use microtype if available
  \usepackage[]{microtype}
  \UseMicrotypeSet[protrusion]{basicmath} % disable protrusion for tt fonts
}{}
\usepackage{setspace}
% Make \paragraph and \subparagraph free-standing
\makeatletter
\ifx\paragraph\undefined\else
  \let\oldparagraph\paragraph
  \renewcommand{\paragraph}{
    \@ifstar
      \xxxParagraphStar
      \xxxParagraphNoStar
  }
  \newcommand{\xxxParagraphStar}[1]{\oldparagraph*{#1}\mbox{}}
  \newcommand{\xxxParagraphNoStar}[1]{\oldparagraph{#1}\mbox{}}
\fi
\ifx\subparagraph\undefined\else
  \let\oldsubparagraph\subparagraph
  \renewcommand{\subparagraph}{
    \@ifstar
      \xxxSubParagraphStar
      \xxxSubParagraphNoStar
  }
  \newcommand{\xxxSubParagraphStar}[1]{\oldsubparagraph*{#1}\mbox{}}
  \newcommand{\xxxSubParagraphNoStar}[1]{\oldsubparagraph{#1}\mbox{}}
\fi
\makeatother

\usepackage{color}
\usepackage{fancyvrb}
\newcommand{\VerbBar}{|}
\newcommand{\VERB}{\Verb[commandchars=\\\{\}]}
\DefineVerbatimEnvironment{Highlighting}{Verbatim}{commandchars=\\\{\}}
% Add ',fontsize=\small' for more characters per line
\usepackage{framed}
\definecolor{shadecolor}{RGB}{241,243,245}
\newenvironment{Shaded}{\begin{snugshade}}{\end{snugshade}}
\newcommand{\AlertTok}[1]{\textcolor[rgb]{0.68,0.00,0.00}{#1}}
\newcommand{\AnnotationTok}[1]{\textcolor[rgb]{0.37,0.37,0.37}{#1}}
\newcommand{\AttributeTok}[1]{\textcolor[rgb]{0.40,0.45,0.13}{#1}}
\newcommand{\BaseNTok}[1]{\textcolor[rgb]{0.68,0.00,0.00}{#1}}
\newcommand{\BuiltInTok}[1]{\textcolor[rgb]{0.00,0.23,0.31}{#1}}
\newcommand{\CharTok}[1]{\textcolor[rgb]{0.13,0.47,0.30}{#1}}
\newcommand{\CommentTok}[1]{\textcolor[rgb]{0.37,0.37,0.37}{#1}}
\newcommand{\CommentVarTok}[1]{\textcolor[rgb]{0.37,0.37,0.37}{\textit{#1}}}
\newcommand{\ConstantTok}[1]{\textcolor[rgb]{0.56,0.35,0.01}{#1}}
\newcommand{\ControlFlowTok}[1]{\textcolor[rgb]{0.00,0.23,0.31}{\textbf{#1}}}
\newcommand{\DataTypeTok}[1]{\textcolor[rgb]{0.68,0.00,0.00}{#1}}
\newcommand{\DecValTok}[1]{\textcolor[rgb]{0.68,0.00,0.00}{#1}}
\newcommand{\DocumentationTok}[1]{\textcolor[rgb]{0.37,0.37,0.37}{\textit{#1}}}
\newcommand{\ErrorTok}[1]{\textcolor[rgb]{0.68,0.00,0.00}{#1}}
\newcommand{\ExtensionTok}[1]{\textcolor[rgb]{0.00,0.23,0.31}{#1}}
\newcommand{\FloatTok}[1]{\textcolor[rgb]{0.68,0.00,0.00}{#1}}
\newcommand{\FunctionTok}[1]{\textcolor[rgb]{0.28,0.35,0.67}{#1}}
\newcommand{\ImportTok}[1]{\textcolor[rgb]{0.00,0.46,0.62}{#1}}
\newcommand{\InformationTok}[1]{\textcolor[rgb]{0.37,0.37,0.37}{#1}}
\newcommand{\KeywordTok}[1]{\textcolor[rgb]{0.00,0.23,0.31}{\textbf{#1}}}
\newcommand{\NormalTok}[1]{\textcolor[rgb]{0.00,0.23,0.31}{#1}}
\newcommand{\OperatorTok}[1]{\textcolor[rgb]{0.37,0.37,0.37}{#1}}
\newcommand{\OtherTok}[1]{\textcolor[rgb]{0.00,0.23,0.31}{#1}}
\newcommand{\PreprocessorTok}[1]{\textcolor[rgb]{0.68,0.00,0.00}{#1}}
\newcommand{\RegionMarkerTok}[1]{\textcolor[rgb]{0.00,0.23,0.31}{#1}}
\newcommand{\SpecialCharTok}[1]{\textcolor[rgb]{0.37,0.37,0.37}{#1}}
\newcommand{\SpecialStringTok}[1]{\textcolor[rgb]{0.13,0.47,0.30}{#1}}
\newcommand{\StringTok}[1]{\textcolor[rgb]{0.13,0.47,0.30}{#1}}
\newcommand{\VariableTok}[1]{\textcolor[rgb]{0.07,0.07,0.07}{#1}}
\newcommand{\VerbatimStringTok}[1]{\textcolor[rgb]{0.13,0.47,0.30}{#1}}
\newcommand{\WarningTok}[1]{\textcolor[rgb]{0.37,0.37,0.37}{\textit{#1}}}

\usepackage{longtable,booktabs,array}
\usepackage{calc} % for calculating minipage widths
% Correct order of tables after \paragraph or \subparagraph
\usepackage{etoolbox}
\makeatletter
\patchcmd\longtable{\par}{\if@noskipsec\mbox{}\fi\par}{}{}
\makeatother
% Allow footnotes in longtable head/foot
\IfFileExists{footnotehyper.sty}{\usepackage{footnotehyper}}{\usepackage{footnote}}
\makesavenoteenv{longtable}
\usepackage{graphicx}
\makeatletter
\newsavebox\pandoc@box
\newcommand*\pandocbounded[1]{% scales image to fit in text height/width
  \sbox\pandoc@box{#1}%
  \Gscale@div\@tempa{\textheight}{\dimexpr\ht\pandoc@box+\dp\pandoc@box\relax}%
  \Gscale@div\@tempb{\linewidth}{\wd\pandoc@box}%
  \ifdim\@tempb\p@<\@tempa\p@\let\@tempa\@tempb\fi% select the smaller of both
  \ifdim\@tempa\p@<\p@\scalebox{\@tempa}{\usebox\pandoc@box}%
  \else\usebox{\pandoc@box}%
  \fi%
}
% Set default figure placement to htbp
\def\fps@figure{htbp}
\makeatother



\ifLuaTeX
\usepackage[bidi=basic,provide=*]{babel}
\else
\usepackage[bidi=default,provide=*]{babel}
\fi
% get rid of language-specific shorthands (see #6817):
\let\LanguageShortHands\languageshorthands
\def\languageshorthands#1{}


\setlength{\emergencystretch}{3em} % prevent overfull lines

\providecommand{\tightlist}{%
  \setlength{\itemsep}{0pt}\setlength{\parskip}{0pt}}



 
\usepackage[style=gost-numeric,backend=biber,langhook=extras,autolang=other*]{biblatex}
\addbibresource{bib/cite.bib}

\usepackage[]{csquotes}

\usepackage[T2A]{fontenc}
\usepackage[utf8]{inputenc}
\usepackage[russian]{babel}

\usepackage{indentfirst}
\setlength{\parindent}{1.25cm}

\usepackage{float}
\floatplacement{figure}{H}

\IfFileExists{plex-otf.sty}{
  \usepackage[
    RM={Scale=0.94},
    SS={Scale=0.94},
    SScon={Scale=0.94},
    TT={Scale=MatchLowercase,FakeStretch=0.9},
    DefaultFeatures={Ligatures=Common}
  ]{plex-otf}
}{%
  \usepackage{libertine}
  \usepackage{juliamono}
}

\usepackage{csquotes}
\KOMAoption{captions}{tableheading}
\makeatletter
\@ifpackageloaded{caption}{}{\usepackage{caption}}
\AtBeginDocument{%
\ifdefined\contentsname
  \renewcommand*\contentsname{Содержание}
\else
  \newcommand\contentsname{Содержание}
\fi
\ifdefined\listfigurename
  \renewcommand*\listfigurename{Список иллюстраций}
\else
  \newcommand\listfigurename{Список иллюстраций}
\fi
\ifdefined\listtablename
  \renewcommand*\listtablename{Список таблиц}
\else
  \newcommand\listtablename{Список таблиц}
\fi
\ifdefined\figurename
  \renewcommand*\figurename{Рисунок}
\else
  \newcommand\figurename{Рисунок}
\fi
\ifdefined\tablename
  \renewcommand*\tablename{Таблица}
\else
  \newcommand\tablename{Таблица}
\fi
}
\@ifpackageloaded{float}{}{\usepackage{float}}
\floatstyle{ruled}
\@ifundefined{c@chapter}{\newfloat{codelisting}{h}{lop}}{\newfloat{codelisting}{h}{lop}[chapter]}
\floatname{codelisting}{Список}
\newcommand*\listoflistings{\listof{codelisting}{Листинги}}
\makeatother
\makeatletter
\makeatother
\makeatletter
\@ifpackageloaded{caption}{}{\usepackage{caption}}
\@ifpackageloaded{subcaption}{}{\usepackage{subcaption}}
\makeatother
\usepackage{bookmark}
\IfFileExists{xurl.sty}{\usepackage{xurl}}{} % add URL line breaks if available
\urlstyle{same}
\hypersetup{
  pdftitle={Лабораторная работа №3},
  pdfauthor={Люпп Софья Романовна},
  pdflang={ru-RU},
  hidelinks,
  pdfcreator={LaTeX via pandoc}}


\title{Лабораторная работа №3}
\usepackage{etoolbox}
\makeatletter
\providecommand{\subtitle}[1]{% add subtitle to \maketitle
  \apptocmd{\@title}{\par {\large #1 \par}}{}{}
}
\makeatother
\subtitle{Агентное моделирование}
\author{Люпп Софья Романовна}
\date{}
\begin{document}
\maketitle

\renewcommand*\contentsname{Содержание}
{
\setcounter{tocdepth}{1}
\tableofcontents
}
\listoffigures
\listoftables

\setstretch{1.5}
\chapter{Цель
работы}\label{ux446ux435ux43bux44c-ux440ux430ux431ux43eux442ux44b}

Ознакомиться с принципами агентного моделирования. Изучить модель
daisyworld. Построить графики daisyworld в соответствии с принципами
агентного моделирования.

\chapter{Задание}\label{ux437ux430ux434ux430ux43dux438ux435}

--- Установить необходимые пакеты. --- Выполнить предложенный код. ---
Преобразовать код в литературный стиль. --- Выполнить документацию в
формате Quarto. --- Выполнить код из jupyter notebook. --- Интегрировать
документацию в формате Quarto в отчёт. --- Добавить в код в литературном
стиле. --- Результирующие файлы выложить на git.

\chapter{Теоретическое
введение}\label{ux442ux435ux43eux440ux435ux442ux438ux447ux435ux441ux43aux43eux435-ux432ux432ux435ux434ux435ux43dux438ux435}

Агентный подход к имитационному моделированию (Agent-Based Modeling,
ABM) --- это метод исследования сложных систем, в котором поведение
системы возникает из взаимодействия множества автономных сущностей,
называемых агентами. Вместо того чтобы описывать систему глобальными
уравнениями, мы моделируем каждую индивидуальную единицу и правила её
поведения, а затем наблюдаем, какие коллективные паттерны появляются
снизу вверх.

Агентное моделирование используется в экологии (популяционная динамика,
модели вроде Daisyworld), экономике (рынки, поведение потребителей),
социологии (распространение мнений, сегрегация), эпидемиологии
(распространение инфекций), транспорте (движение пешеходов и
автомобилей), политологии и многих других областях.

Модель Daisyworld, предложенная Джеймсом Лавлоком и Эндрю Уотсоном,
иллюстрирует гипотезу Геи, которая рассматривает планету как единую,
саморегулирующуюся систему, включающую как живые, так и неживые части. В
мире маргариток произрастают чёрные и белые маргаритки. Их альбедо
различается: чёрные маргаритки поглощают свет и тепло, нагревая
окружающую среду; белые маргаритки делают наоборот. Маргаритки могут
размножаться только в определённом температурном диапазоне, а это
значит, что слишком много (или слишком мало) тепла от солнца и/или
окружающей среды в конечном итоге остановит их размножение.

В модели Daisyworld агентами являются чёрные и белые маргаритки. Они
живут на клеточной сетке (среда). Свойства агентов: вид и возраст.
Правила: --- Маргаритки изменяют локальную температуру за счёт разного
альбедо. --- Температура влияет на вероятность размножения (чем ближе к
оптимуму, тем выше шанс заселить соседнюю пустую клетку). --- Агенты
стареют и умирают после определённого возраста. --- Среда (температура)
диффундирует между клетками.

\chapter{Выполнение лабораторной
работы}\label{ux432ux44bux43fux43eux43bux43dux435ux43dux438ux435-ux43bux430ux431ux43eux440ux430ux442ux43eux440ux43dux43eux439-ux440ux430ux431ux43eux442ux44b}

Начинаю выполнение лабораторной работы с установки пакета Agents, и
создания нового файла daisyworld.jl в scr и scripts (рис.~\ref{fig-001})
(рис.~\ref{fig-002}) (рис.~\ref{fig-003}).

\begin{figure}

\centering{

\includegraphics[width=0.7\linewidth,height=\textheight,keepaspectratio]{image/1.png}

}

\caption{\label{fig-001}Создание файла daisyworld.jl}

\end{figure}%

\begin{figure}

\centering{

\includegraphics[width=0.7\linewidth,height=\textheight,keepaspectratio]{image/2.png}

}

\caption{\label{fig-002}Установка пакета Agents}

\end{figure}%

\begin{figure}

\centering{

\includegraphics[width=0.7\linewidth,height=\textheight,keepaspectratio]{image/3.png}

}

\caption{\label{fig-003}Установка пакета Agents}

\end{figure}%

Заполняю код для файла daisyworld.jl, который определяет тип агента и
функции шага модели (рис.~\ref{fig-004}).

\begin{figure}

\centering{

\includegraphics[width=0.7\linewidth,height=\textheight,keepaspectratio]{image/4.png}

}

\caption{\label{fig-004}daisyworld}

\end{figure}%

Запускаю код через Julia, получая в plots графики - тепловые карты,
которые будут представлять собой температуру поверхности модели.
Маргаритки будут отображаться черно-белыми в соответствии с их видом
(рис.~\ref{fig-005}) (рис.~\ref{fig-006}) (рис.~\ref{fig-007})
(рис.~\ref{fig-008}).

\begin{figure}

\centering{

\includegraphics[width=0.7\linewidth,height=\textheight,keepaspectratio]{image/5.png}

}

\caption{\label{fig-005}daisyworld plots}

\end{figure}%

\begin{figure}

\centering{

\includegraphics[width=0.7\linewidth,height=\textheight,keepaspectratio]{image/6.png}

}

\caption{\label{fig-006}Тепловая карта daisyworld}

\end{figure}%

\begin{figure}

\centering{

\includegraphics[width=0.7\linewidth,height=\textheight,keepaspectratio]{image/7.png}

}

\caption{\label{fig-007}Тепловая карта daisyworld}

\end{figure}%

\begin{figure}

\centering{

\includegraphics[width=0.7\linewidth,height=\textheight,keepaspectratio]{image/8.png}

}

\caption{\label{fig-008}Тепловая карта daisyworld}

\end{figure}%

Создадим визуализацию не отдельных моментов, а видео эволюции модели при
помощи нового файла daisyworld-animate.jl (рис.~\ref{fig-009}).

\begin{figure}

\centering{

\includegraphics[width=0.7\linewidth,height=\textheight,keepaspectratio]{image/9.png}

}

\caption{\label{fig-009}Анимация эволюции модели}

\end{figure}%

На выходе получаем видео формата mp4 эволюции модели
(рис.~\ref{fig-010}) (рис.~\ref{fig-011}).

\begin{figure}

\centering{

\includegraphics[width=0.7\linewidth,height=\textheight,keepaspectratio]{image/10.png}

}

\caption{\label{fig-010}Видео эволюции модели}

\end{figure}%

\begin{figure}

\centering{

\includegraphics[width=0.7\linewidth,height=\textheight,keepaspectratio]{image/11.png}

}

\caption{\label{fig-011}Анимация эволюции модели}

\end{figure}%

Далее построим график daisyworld-count.jl изменения числа маргариток в
зависимости от модельного времени (рис.~\ref{fig-012})
(рис.~\ref{fig-013}) (рис.~\ref{fig-014}).

\begin{figure}

\centering{

\includegraphics[width=0.7\linewidth,height=\textheight,keepaspectratio]{image/12.png}

}

\caption{\label{fig-012}daisyworld-count в папке}

\end{figure}%

\begin{figure}

\centering{

\includegraphics[width=0.7\linewidth,height=\textheight,keepaspectratio]{image/13.png}

}

\caption{\label{fig-013}График модели daisyworld-count}

\end{figure}%

\begin{figure}

\centering{

\includegraphics[width=0.7\linewidth,height=\textheight,keepaspectratio]{image/14.png}

}

\caption{\label{fig-014}daisyworld-count}

\end{figure}%

После построим комплексный график изменения числа маргариток,
температуры, альбедо в зависимости от модельного времени
(рис.~\ref{fig-015}) (рис.~\ref{fig-016}).

\begin{figure}

\centering{

\includegraphics[width=0.7\linewidth,height=\textheight,keepaspectratio]{image/15.png}

}

\caption{\label{fig-015}График модели daisyworld-luminosity}

\end{figure}%

\begin{figure}

\centering{

\includegraphics[width=0.7\linewidth,height=\textheight,keepaspectratio]{image/16.png}

}

\caption{\label{fig-016}daisyworld-luminosity}

\end{figure}%

Далее делаю производные форматы с помощи Julia и файла tangle.jl всех
новых файлов (рис.~\ref{fig-017}) (рис.~\ref{fig-018})
(рис.~\ref{fig-019}).

\begin{figure}

\centering{

\includegraphics[width=0.7\linewidth,height=\textheight,keepaspectratio]{image/17.png}

}

\caption{\label{fig-017}Производные форматы модели daisyworld.jl}

\end{figure}%

\begin{figure}

\centering{

\includegraphics[width=0.7\linewidth,height=\textheight,keepaspectratio]{image/18.png}

}

\caption{\label{fig-018}Jupyter notebook daisyworld.jl}

\end{figure}%

\begin{figure}

\centering{


\includegraphics[width=0.7\linewidth,height=\textheight,keepaspectratio]{image/19.png}

}

\caption{\label{fig-019}Jupyter notebook daisyworld}

\end{figure}%

Создаю производные форматы для файла daisyworld-count.jl
(рис.~\ref{fig-020}) (рис.~\ref{fig-021}).

\begin{figure}

\centering{

\includegraphics[width=0.7\linewidth,height=\textheight,keepaspectratio]{image/20.png}

}

\caption{\label{fig-020}daisyworld-count.jl}

\end{figure}%

\begin{figure}

\centering{

\includegraphics[width=0.7\linewidth,height=\textheight,keepaspectratio]{image/21.png}

}

\caption{\label{fig-021}Jupyter notebook daisyworld-count.jl}

\end{figure}%

Для расширения базовой визуализации за счёт параметров создам файл
scripts daisyworld\_\_param.jl (рис.~\ref{fig-022}) (рис.~\ref{fig-021})
(рис.~\ref{fig-022}) (рис.~\ref{fig-023}) (рис.~\ref{fig-024})
(рис.~\ref{fig-025}) (рис.~\ref{fig-026}) (рис.~\ref{fig-027})
(рис.~\ref{fig-028}) (рис.~\ref{fig-029}) (рис.~\ref{fig-030})
(рис.~\ref{fig-031}) (рис.~\ref{fig-032}) (рис.~\ref{fig-033})
(рис.~\ref{fig-034}).

\begin{figure}

\centering{

\includegraphics[width=0.7\linewidth,height=\textheight,keepaspectratio]{image/22.png}

}

\caption{\label{fig-022}Создание файла daisyworld\_\_param.jl}

\end{figure}%

\begin{figure}

\centering{

\includegraphics[width=0.7\linewidth,height=\textheight,keepaspectratio]{image/23.png}

}

\caption{\label{fig-023}daisyworld\_\_param.jl}

\end{figure}%

\begin{figure}

\centering{

\includegraphics[width=0.7\linewidth,height=\textheight,keepaspectratio]{image/24.png}

}

\caption{\label{fig-024}daisyworld\_\_param.jl}

\end{figure}%

\begin{figure}

\centering{

\includegraphics[width=0.7\linewidth,height=\textheight,keepaspectratio]{image/25.png}

}

\caption{\label{fig-025}daisyworld\_\_param.jl}

\end{figure}%

\begin{figure}

\centering{

\includegraphics[width=0.7\linewidth,height=\textheight,keepaspectratio]{image/26.png}

}

\caption{\label{fig-026}daisyworld\_\_param.jl}

\end{figure}%

\begin{figure}

\centering{

\includegraphics[width=0.7\linewidth,height=\textheight,keepaspectratio]{image/27.png}

}

\caption{\label{fig-027}daisyworld\_\_param.jl}

\end{figure}%

\begin{figure}

\centering{

\includegraphics[width=0.7\linewidth,height=\textheight,keepaspectratio]{image/28.png}

}

\caption{\label{fig-028}daisyworld\_\_param.jl}

\end{figure}%

\begin{figure}

\centering{

\includegraphics[width=0.7\linewidth,height=\textheight,keepaspectratio]{image/29.png}

}

\caption{\label{fig-029}daisyworld\_\_param.jl}

\end{figure}%

\begin{figure}

\centering{

\includegraphics[width=0.7\linewidth,height=\textheight,keepaspectratio]{image/30.png}

}

\caption{\label{fig-030}daisyworld\_\_param.jl}

\end{figure}%

\begin{figure}

\centering{

\includegraphics[width=0.7\linewidth,height=\textheight,keepaspectratio]{image/31.png}

}

\caption{\label{fig-031}daisyworld\_\_param.jl}

\end{figure}%

\begin{figure}

\centering{

\includegraphics[width=0.7\linewidth,height=\textheight,keepaspectratio]{image/32.png}

}

\caption{\label{fig-032}daisyworld\_\_param.jl}

\end{figure}%

\begin{figure}

\centering{

\includegraphics[width=0.7\linewidth,height=\textheight,keepaspectratio]{image/33.png}

}

\caption{\label{fig-033}daisyworld\_\_param.jl}

\end{figure}%

\begin{figure}

\centering{

\includegraphics[width=0.7\linewidth,height=\textheight,keepaspectratio]{image/34.png}

}

\caption{\label{fig-034}daisyworld\_\_param.jl}

\end{figure}%

Создаю производные форматы файла daisyworld\_\_param.jl

\begin{figure}

\centering{

\includegraphics[width=0.7\linewidth,height=\textheight,keepaspectratio]{image/35.png}

}

\caption{\label{fig-035}Производные форматы daisyworld\_\_param.jl}

\end{figure}%

Построю график изменения числа маргариток в зависимости от модельного
времени с разными параметрами модели. Для этого создаю файл
scripts/daisyworld-count\_\_param.jl. Запускаю файл через Julia
(рис.~\ref{fig-036}) (рис.~\ref{fig-037}) (рис.~\ref{fig-038})
(рис.~\ref{fig-039}) (рис.~\ref{fig-040}).

\begin{figure}

\centering{

\includegraphics[width=0.7\linewidth,height=\textheight,keepaspectratio]{image/36.png}

}

\caption{\label{fig-036}daisyworld-count\_\_param.jl}

\end{figure}%

\begin{figure}

\centering{

\includegraphics[width=0.7\linewidth,height=\textheight,keepaspectratio]{image/37.png}

}

\caption{\label{fig-037}daisyworld-count\_\_param.jl}

\end{figure}%

\begin{figure}

\centering{

\includegraphics[width=0.7\linewidth,height=\textheight,keepaspectratio]{image/38.png}

}

\caption{\label{fig-038}daisyworld-count\_\_param.jl}

\end{figure}%

\begin{figure}

\centering{

\includegraphics[width=0.7\linewidth,height=\textheight,keepaspectratio]{image/39.png}

}

\caption{\label{fig-039}daisyworld-count\_\_param.jl}

\end{figure}%

\begin{figure}

\centering{

\includegraphics[width=0.7\linewidth,height=\textheight,keepaspectratio]{image/40.png}

}

\caption{\label{fig-040}Создание файла daisyworld-count\_\_param.jl}

\end{figure}%

Построю комплексный график изменения числа маргариток, температуры,
альбедо в зависимости от модельного времени с разными параметрами модели
(рис.~\ref{fig-043}) (рис.~\ref{fig-044}) (рис.~\ref{fig-045})
(рис.~\ref{fig-046}). Создам файл
scripts/daisyworld-luminosity\_\_param.jl (рис.~\ref{fig-041}). Далее
также делаю производные форматы (рис.~\ref{fig-042})

\begin{figure}

\centering{

\includegraphics[width=0.7\linewidth,height=\textheight,keepaspectratio]{image/41.png}

}

\caption{\label{fig-041}Создание файла
daisyworld-luminosity\_\_param.jl}

\end{figure}%

\begin{figure}

\centering{

\includegraphics[width=0.7\linewidth,height=\textheight,keepaspectratio]{image/42.png}

}

\caption{\label{fig-042}Создание производных форматов}

\end{figure}%

\begin{figure}

\centering{

\includegraphics[width=0.7\linewidth,height=\textheight,keepaspectratio]{image/43.png}

}

\caption{\label{fig-043}daisyworld-luminosity\_\_param.jl}

\end{figure}%

\begin{figure}

\centering{

\includegraphics[width=0.7\linewidth,height=\textheight,keepaspectratio]{image/44.png}

}

\caption{\label{fig-044}daisyworld-luminosity\_\_param.jl}

\end{figure}%

\begin{figure}

\centering{

\includegraphics[width=0.7\linewidth,height=\textheight,keepaspectratio]{image/45.png}

}

\caption{\label{fig-045}daisyworld-luminosity\_\_param.jl}

\end{figure}%

\begin{figure}

\centering{

\includegraphics[width=0.7\linewidth,height=\textheight,keepaspectratio]{image/46.png}

}

\caption{\label{fig-046}daisyworld-luminosity\_\_param.jl}

\end{figure}%

Подключаю Julia в файле quarto лабораторной работы (рис.~\ref{fig-047}).

\begin{figure}

\centering{

\includegraphics[width=0.7\linewidth,height=\textheight,keepaspectratio]{image/47.png}

}

\caption{\label{fig-047}Подключение Julia}

\end{figure}%

\begin{Shaded}
\begin{Highlighting}[]
\ImportTok{using} \BuiltInTok{DrWatson}
\PreprocessorTok{@quickactivate} \StringTok{"project"}
\ImportTok{using} \BuiltInTok{Agents}
\ImportTok{using} \BuiltInTok{DataFrames}
\ImportTok{using} \BuiltInTok{Plots}
\FunctionTok{include}\NormalTok{(}\StringTok{"/home/srluipp/project/src/daisyworld.jl"}\NormalTok{)}
\ImportTok{using} \BuiltInTok{CairoMakie}
\NormalTok{model }\OperatorTok{=} \FunctionTok{daisyworld}\NormalTok{()}
\FunctionTok{daisycolor}\NormalTok{(a}\OperatorTok{::}\DataTypeTok{Daisy}\NormalTok{) }\OperatorTok{=}\NormalTok{ a.breed}
\NormalTok{plotkwargs }\OperatorTok{=}\NormalTok{ (}
\NormalTok{agent\_color}\OperatorTok{=}\NormalTok{daisycolor, agent\_size }\OperatorTok{=} \FloatTok{20}\NormalTok{, agent\_marker }\OperatorTok{=} \CharTok{\textquotesingle{}✿\textquotesingle{}}\NormalTok{,}
\NormalTok{heatarray }\OperatorTok{=} \OperatorTok{:}\NormalTok{temperature,}
\NormalTok{heatkwargs }\OperatorTok{=}\NormalTok{ (colorrange }\OperatorTok{=}\NormalTok{ (}\OperatorTok{{-}}\FloatTok{20}\NormalTok{, }\FloatTok{60}\NormalTok{),),}
\NormalTok{)}
\NormalTok{plt1, \_ }\OperatorTok{=} \FunctionTok{abmplot}\NormalTok{(model; plotkwargs}\OperatorTok{...}\NormalTok{)}
\FunctionTok{step!}\NormalTok{(model, }\FloatTok{5}\NormalTok{)}
\NormalTok{plt2, \_ }\OperatorTok{=} \FunctionTok{abmplot}\NormalTok{(model; heatarray }\OperatorTok{=}\NormalTok{ model.temperature, plotkwargs}\OperatorTok{...}\NormalTok{)}
\FunctionTok{step!}\NormalTok{(model, }\FloatTok{40}\NormalTok{)}
\NormalTok{plt3, \_ }\OperatorTok{=} \FunctionTok{abmplot}\NormalTok{(model; heatarray }\OperatorTok{=}\NormalTok{ model.temperature, plotkwargs}\OperatorTok{...}\NormalTok{)}
\FunctionTok{save}\NormalTok{(}\FunctionTok{plotsdir}\NormalTok{(}\StringTok{"daisy\_step001.png"}\NormalTok{), plt1)}
\FunctionTok{save}\NormalTok{(}\FunctionTok{plotsdir}\NormalTok{(}\StringTok{"daisy\_step005.png"}\NormalTok{), plt2)}
\FunctionTok{save}\NormalTok{(}\FunctionTok{plotsdir}\NormalTok{(}\StringTok{"daisy\_step040.png"}\NormalTok{), plt3)}
\end{Highlighting}
\end{Shaded}

\begin{verbatim}
┌ Warning: Found `resolution` in the theme when creating a `Scene`. The `resolution` keyword for `Scene`s and `Figure`s has been deprecated. Use `Figure(; size = ...` or `Scene(; size = ...)` instead, which better reflects that this is a unitless size and not a pixel resolution. The key could also come from `set_theme!` calls or related theming functions.
└ @ Makie ~/.julia/packages/Makie/kJl0u/src/scenes.jl:264
┌ Warning: Found `resolution` in the theme when creating a `Scene`. The `resolution` keyword for `Scene`s and `Figure`s has been deprecated. Use `Figure(; size = ...` or `Scene(; size = ...)` instead, which better reflects that this is a unitless size and not a pixel resolution. The key could also come from `set_theme!` calls or related theming functions.
└ @ Makie ~/.julia/packages/Makie/kJl0u/src/scenes.jl:264
┌ Warning: Found `resolution` in the theme when creating a `Scene`. The `resolution` keyword for `Scene`s and `Figure`s has been deprecated. Use `Figure(; size = ...` or `Scene(; size = ...)` instead, which better reflects that this is a unitless size and not a pixel resolution. The key could also come from `set_theme!` calls or related theming functions.
└ @ Makie ~/.julia/packages/Makie/kJl0u/src/scenes.jl:264
\end{verbatim}

\begin{Shaded}
\begin{Highlighting}[]
\ImportTok{using} \BuiltInTok{DrWatson}
\PreprocessorTok{@quickactivate} \StringTok{"project"}
\ImportTok{using} \BuiltInTok{Agents}
\ImportTok{using} \BuiltInTok{DataFrames}
\ImportTok{using} \BuiltInTok{Plots}
\FunctionTok{include}\NormalTok{(}\StringTok{"/home/srluipp/project/src/daisyworld.jl"}\NormalTok{)}
\ImportTok{using} \BuiltInTok{CairoMakie}
\FunctionTok{black}\NormalTok{(a) }\OperatorTok{=}\NormalTok{ a.breed }\OperatorTok{==} \OperatorTok{:}\NormalTok{black}
\FunctionTok{white}\NormalTok{(a) }\OperatorTok{=}\NormalTok{ a.breed }\OperatorTok{==} \OperatorTok{:}\NormalTok{white}
\NormalTok{adata }\OperatorTok{=}\NormalTok{ [(black, count), (white, count)]}
\NormalTok{model }\OperatorTok{=} \FunctionTok{daisyworld}\NormalTok{(; solar\_luminosity }\OperatorTok{=} \FloatTok{1.0}\NormalTok{)}
\NormalTok{agent\_df, model\_df }\OperatorTok{=} \FunctionTok{run!}\NormalTok{(model, }\FloatTok{1000}\NormalTok{; adata)}
\NormalTok{figure }\OperatorTok{=} \FunctionTok{Figure}\NormalTok{(size }\OperatorTok{=}\NormalTok{ (}\FloatTok{600}\NormalTok{, }\FloatTok{400}\NormalTok{));}
\NormalTok{ax }\OperatorTok{=}\NormalTok{ figure[}\FloatTok{1}\NormalTok{, }\FloatTok{1}\NormalTok{] }\OperatorTok{=} \FunctionTok{Axis}\NormalTok{(figure, xlabel }\OperatorTok{=} \StringTok{"tick"}\NormalTok{, ylabel }\OperatorTok{=} \StringTok{"daisy count"}\NormalTok{)}
\NormalTok{blackl }\OperatorTok{=} \FunctionTok{lines!}\NormalTok{(ax, agent\_df[!, }\OperatorTok{:}\NormalTok{time], agent\_df[!, }\OperatorTok{:}\NormalTok{count\_black], color }\OperatorTok{=} \OperatorTok{:}\NormalTok{black)}
\NormalTok{whitel }\OperatorTok{=} \FunctionTok{lines!}\NormalTok{(ax, agent\_df[!, }\OperatorTok{:}\NormalTok{time], agent\_df[!, }\OperatorTok{:}\NormalTok{count\_white], color }\OperatorTok{=} \OperatorTok{:}\NormalTok{orange)}
\FunctionTok{Legend}\NormalTok{(figure[}\FloatTok{1}\NormalTok{, }\FloatTok{2}\NormalTok{], [blackl, whitel], [}\StringTok{"black"}\NormalTok{, }\StringTok{"white"}\NormalTok{], labelsize }\OperatorTok{=} \FloatTok{12}\NormalTok{)}
\end{Highlighting}
\end{Shaded}

\begin{verbatim}
┌ Warning: Found `resolution` in the theme when creating a `Scene`. The `resolution` keyword for `Scene`s and `Figure`s has been deprecated. Use `Figure(; size = ...` or `Scene(; size = ...)` instead, which better reflects that this is a unitless size and not a pixel resolution. The key could also come from `set_theme!` calls or related theming functions.
└ @ Makie ~/.julia/packages/Makie/kJl0u/src/scenes.jl:264
\end{verbatim}

\begin{verbatim}
Legend()
\end{verbatim}

figure

\begin{Shaded}
\begin{Highlighting}[]
\FunctionTok{save}\NormalTok{(}\FunctionTok{plotsdir}\NormalTok{(}\StringTok{"daisy\_count.png"}\NormalTok{), figure)}
\end{Highlighting}
\end{Shaded}

\begin{Shaded}
\begin{Highlighting}[]
\ImportTok{using} \BuiltInTok{DrWatson}
\PreprocessorTok{@quickactivate} \StringTok{"project"}
\ImportTok{using} \BuiltInTok{Agents}
\ImportTok{using} \BuiltInTok{DataFrames}
\ImportTok{using} \BuiltInTok{Plots}
\ImportTok{using} \BuiltInTok{CairoMakie}
\FunctionTok{include}\NormalTok{(}\StringTok{"/home/srluipp/project/src/daisyworld.jl"}\NormalTok{)}
\FunctionTok{black}\NormalTok{(a) }\OperatorTok{=}\NormalTok{ a.breed }\OperatorTok{==} \OperatorTok{:}\NormalTok{black}
\FunctionTok{white}\NormalTok{(a) }\OperatorTok{=}\NormalTok{ a.breed }\OperatorTok{==} \OperatorTok{:}\NormalTok{white}
\NormalTok{adata }\OperatorTok{=}\NormalTok{ [(black, count), (white, count)]}
\CommentTok{\# Параметры эксперимента}
\NormalTok{param\_dict }\OperatorTok{=} \FunctionTok{Dict}\NormalTok{(}
\OperatorTok{:}\NormalTok{griddims }\OperatorTok{=\textgreater{}}\NormalTok{ (}\FloatTok{30}\NormalTok{, }\FloatTok{30}\NormalTok{),}
\OperatorTok{:}\NormalTok{max\_age }\OperatorTok{=\textgreater{}}\NormalTok{ [}\FloatTok{25}\NormalTok{, }\FloatTok{40}\NormalTok{],}
\OperatorTok{:}\NormalTok{init\_white }\OperatorTok{=\textgreater{}}\NormalTok{ [}\FloatTok{0.2}\NormalTok{, }\FloatTok{0.8}\NormalTok{],}
\OperatorTok{:}\NormalTok{init\_black }\OperatorTok{=\textgreater{}} \FloatTok{0.2}\NormalTok{,}
\OperatorTok{:}\NormalTok{albedo\_white }\OperatorTok{=\textgreater{}} \FloatTok{0.75}\NormalTok{,}
\OperatorTok{:}\NormalTok{albedo\_black }\OperatorTok{=\textgreater{}} \FloatTok{0.25}\NormalTok{,}
\OperatorTok{:}\NormalTok{surface\_albedo }\OperatorTok{=\textgreater{}} \FloatTok{0.4}\NormalTok{,}
\OperatorTok{:}\NormalTok{solar\_change }\OperatorTok{=\textgreater{}} \FloatTok{0.005}\NormalTok{,}
\OperatorTok{:}\NormalTok{solar\_luminosity }\OperatorTok{=\textgreater{}} \FloatTok{1.0}\NormalTok{,}
\OperatorTok{:}\NormalTok{scenario }\OperatorTok{=\textgreater{}} \OperatorTok{:}\NormalTok{default,}
\OperatorTok{:}\NormalTok{seed }\OperatorTok{=\textgreater{}} \FloatTok{165}\NormalTok{,}
\NormalTok{)}
\CommentTok{\# Создаём список всех комбинаций}
\NormalTok{params\_list }\OperatorTok{=} \FunctionTok{dict\_list}\NormalTok{(param\_dict)}
\ControlFlowTok{for}\NormalTok{ params }\KeywordTok{in}\NormalTok{ params\_list}
\NormalTok{model }\OperatorTok{=} \FunctionTok{daisyworld}\NormalTok{(;params}\OperatorTok{...}\NormalTok{)}
\NormalTok{agent\_df, model\_df }\OperatorTok{=} \FunctionTok{run!}\NormalTok{(model, }\FloatTok{1000}\NormalTok{; adata)}
\NormalTok{figure }\OperatorTok{=} \FunctionTok{Figure}\NormalTok{(size }\OperatorTok{=}\NormalTok{ (}\FloatTok{600}\NormalTok{, }\FloatTok{400}\NormalTok{));}
\NormalTok{ax }\OperatorTok{=}\NormalTok{ figure[}\FloatTok{1}\NormalTok{, }\FloatTok{1}\NormalTok{] }\OperatorTok{=} \FunctionTok{Axis}\NormalTok{(figure, xlabel }\OperatorTok{=} \StringTok{"tick"}\NormalTok{, ylabel }\OperatorTok{=} \StringTok{"daisy count"}\NormalTok{)}
\NormalTok{blackl }\OperatorTok{=} \FunctionTok{lines!}\NormalTok{(ax, agent\_df[!, }\OperatorTok{:}\NormalTok{time], agent\_df[!, }\OperatorTok{:}\NormalTok{count\_black], color }\OperatorTok{=} \OperatorTok{:}\NormalTok{black)}
\NormalTok{whitel }\OperatorTok{=} \FunctionTok{lines!}\NormalTok{(ax, agent\_df[!, }\OperatorTok{:}\NormalTok{time], agent\_df[!, }\OperatorTok{:}\NormalTok{count\_white], color }\OperatorTok{=} \OperatorTok{:}\NormalTok{orange)}
\FunctionTok{Legend}\NormalTok{(figure[}\FloatTok{1}\NormalTok{, }\FloatTok{2}\NormalTok{], [blackl, whitel], [}\StringTok{"black"}\NormalTok{, }\StringTok{"white"}\NormalTok{], labelsize }\OperatorTok{=} \FloatTok{12}\NormalTok{)}
\NormalTok{plt\_name }\OperatorTok{=} \FunctionTok{savename}\NormalTok{(}\StringTok{"daisy{-}count"}\NormalTok{,params) }\OperatorTok{*} \StringTok{".png"}
\FunctionTok{save}\NormalTok{(}\FunctionTok{plotsdir}\NormalTok{(plt\_name), figure)}
\ControlFlowTok{end}
\end{Highlighting}
\end{Shaded}

\begin{verbatim}
┌ Warning: Found `resolution` in the theme when creating a `Scene`. The `resolution` keyword for `Scene`s and `Figure`s has been deprecated. Use `Figure(; size = ...` or `Scene(; size = ...)` instead, which better reflects that this is a unitless size and not a pixel resolution. The key could also come from `set_theme!` calls or related theming functions.
└ @ Makie ~/.julia/packages/Makie/kJl0u/src/scenes.jl:264
┌ Warning: Found `resolution` in the theme when creating a `Scene`. The `resolution` keyword for `Scene`s and `Figure`s has been deprecated. Use `Figure(; size = ...` or `Scene(; size = ...)` instead, which better reflects that this is a unitless size and not a pixel resolution. The key could also come from `set_theme!` calls or related theming functions.
└ @ Makie ~/.julia/packages/Makie/kJl0u/src/scenes.jl:264
┌ Warning: Found `resolution` in the theme when creating a `Scene`. The `resolution` keyword for `Scene`s and `Figure`s has been deprecated. Use `Figure(; size = ...` or `Scene(; size = ...)` instead, which better reflects that this is a unitless size and not a pixel resolution. The key could also come from `set_theme!` calls or related theming functions.
└ @ Makie ~/.julia/packages/Makie/kJl0u/src/scenes.jl:264
┌ Warning: Found `resolution` in the theme when creating a `Scene`. The `resolution` keyword for `Scene`s and `Figure`s has been deprecated. Use `Figure(; size = ...` or `Scene(; size = ...)` instead, which better reflects that this is a unitless size and not a pixel resolution. The key could also come from `set_theme!` calls or related theming functions.
└ @ Makie ~/.julia/packages/Makie/kJl0u/src/scenes.jl:264
\end{verbatim}

\begin{Shaded}
\begin{Highlighting}[]
\ImportTok{using} \BuiltInTok{DrWatson}
\PreprocessorTok{@quickactivate} \StringTok{"project"}
\ImportTok{using} \BuiltInTok{Agents}
\ImportTok{using} \BuiltInTok{DataFrames}
\ImportTok{using} \BuiltInTok{Plots}
\FunctionTok{include}\NormalTok{(}\StringTok{"/home/srluipp/project/src/daisyworld.jl"}\NormalTok{)}
\ImportTok{using} \BuiltInTok{CairoMakie}
\FunctionTok{black}\NormalTok{(a) }\OperatorTok{=}\NormalTok{ a.breed }\OperatorTok{==} \OperatorTok{:}\NormalTok{black}
\FunctionTok{white}\NormalTok{(a) }\OperatorTok{=}\NormalTok{ a.breed }\OperatorTok{==} \OperatorTok{:}\NormalTok{white}
\NormalTok{adata }\OperatorTok{=}\NormalTok{ [(black, count), (white, count)]}
\NormalTok{model }\OperatorTok{=} \FunctionTok{daisyworld}\NormalTok{(solar\_luminosity }\OperatorTok{=} \FloatTok{1.0}\NormalTok{, scenario }\OperatorTok{=} \OperatorTok{:}\NormalTok{ramp)}
\FunctionTok{temperature}\NormalTok{(model) }\OperatorTok{=}\NormalTok{ StatsBase.}\FunctionTok{mean}\NormalTok{(model.temperature)}
\NormalTok{mdata }\OperatorTok{=}\NormalTok{ [temperature, }\OperatorTok{:}\NormalTok{solar\_luminosity]}
\NormalTok{agent\_df, model\_df }\OperatorTok{=} \FunctionTok{run!}\NormalTok{(model, }\FloatTok{1000}\NormalTok{; adata }\OperatorTok{=}\NormalTok{ adata, mdata }\OperatorTok{=}\NormalTok{ mdata)}
\NormalTok{figure }\OperatorTok{=}\NormalTok{ CairoMakie.}\FunctionTok{Figure}\NormalTok{(size }\OperatorTok{=}\NormalTok{ (}\FloatTok{600}\NormalTok{, }\FloatTok{600}\NormalTok{));}
\NormalTok{ax1 }\OperatorTok{=}\NormalTok{ figure[}\FloatTok{1}\NormalTok{, }\FloatTok{1}\NormalTok{] }\OperatorTok{=} \FunctionTok{Axis}\NormalTok{(figure, ylabel }\OperatorTok{=} \StringTok{"daisy count"}\NormalTok{)}
\NormalTok{blackl }\OperatorTok{=} \FunctionTok{lines!}\NormalTok{(ax1, agent\_df[!, }\OperatorTok{:}\NormalTok{time], agent\_df[!, }\OperatorTok{:}\NormalTok{count\_black], color }\OperatorTok{=} \OperatorTok{:}\NormalTok{red)}
\NormalTok{whitel }\OperatorTok{=} \FunctionTok{lines!}\NormalTok{(ax1, agent\_df[!, }\OperatorTok{:}\NormalTok{time], agent\_df[!, }\OperatorTok{:}\NormalTok{count\_white], color }\OperatorTok{=} \OperatorTok{:}\NormalTok{blue)}
\NormalTok{figure[}\FloatTok{1}\NormalTok{, }\FloatTok{2}\NormalTok{] }\OperatorTok{=} \FunctionTok{Legend}\NormalTok{(figure, [blackl, whitel], [}\StringTok{"black"}\NormalTok{, }\StringTok{"white"}\NormalTok{])}
\NormalTok{ax2 }\OperatorTok{=}\NormalTok{ figure[}\FloatTok{2}\NormalTok{, }\FloatTok{1}\NormalTok{] }\OperatorTok{=} \FunctionTok{Axis}\NormalTok{(figure, ylabel }\OperatorTok{=} \StringTok{"temperature"}\NormalTok{)}
\NormalTok{ax3 }\OperatorTok{=}\NormalTok{ figure[}\FloatTok{3}\NormalTok{, }\FloatTok{1}\NormalTok{] }\OperatorTok{=} \FunctionTok{Axis}\NormalTok{(figure, xlabel }\OperatorTok{=} \StringTok{"tick"}\NormalTok{, ylabel }\OperatorTok{=} \StringTok{"luminosity"}\NormalTok{)}
\FunctionTok{lines!}\NormalTok{(ax2, model\_df[!, }\OperatorTok{:}\NormalTok{time], model\_df[!, }\OperatorTok{:}\NormalTok{temperature], color }\OperatorTok{=} \OperatorTok{:}\NormalTok{red)}
\FunctionTok{lines!}\NormalTok{(ax3, model\_df[!, }\OperatorTok{:}\NormalTok{time], model\_df[!, }\OperatorTok{:}\NormalTok{solar\_luminosity], color }\OperatorTok{=} \OperatorTok{:}\NormalTok{red)}
\ControlFlowTok{for}\NormalTok{ ax }\KeywordTok{in}\NormalTok{ (ax1, ax2); ax.xticklabelsvisible }\OperatorTok{=} \ConstantTok{false}\NormalTok{; }\ControlFlowTok{end}
\NormalTok{figure}
\FunctionTok{save}\NormalTok{(}\FunctionTok{plotsdir}\NormalTok{(}\StringTok{"daisy\_luminosity.png"}\NormalTok{), figure)}
\end{Highlighting}
\end{Shaded}

\begin{verbatim}
┌ Warning: Found `resolution` in the theme when creating a `Scene`. The `resolution` keyword for `Scene`s and `Figure`s has been deprecated. Use `Figure(; size = ...` or `Scene(; size = ...)` instead, which better reflects that this is a unitless size and not a pixel resolution. The key could also come from `set_theme!` calls or related theming functions.
└ @ Makie ~/.julia/packages/Makie/kJl0u/src/scenes.jl:264
\end{verbatim}

\begin{Shaded}
\begin{Highlighting}[]
\ImportTok{using} \BuiltInTok{DrWatson}
\PreprocessorTok{@quickactivate} \StringTok{"project"}
\ImportTok{using} \BuiltInTok{Agents}
\ImportTok{using} \BuiltInTok{DataFrames}
\ImportTok{using} \BuiltInTok{Plots}
\ImportTok{using} \BuiltInTok{CairoMakie}
\FunctionTok{include}\NormalTok{(}\StringTok{"/home/srluipp/project/src/daisyworld.jl"}\NormalTok{)}
\FunctionTok{black}\NormalTok{(a) }\OperatorTok{=}\NormalTok{ a.breed }\OperatorTok{==} \OperatorTok{:}\NormalTok{black}
\FunctionTok{white}\NormalTok{(a) }\OperatorTok{=}\NormalTok{ a.breed }\OperatorTok{==} \OperatorTok{:}\NormalTok{white}
\NormalTok{adata }\OperatorTok{=}\NormalTok{ [(black, count), (white, count)]}
\CommentTok{\# Параметры эксперимента}
\NormalTok{param\_dict }\OperatorTok{=} \FunctionTok{Dict}\NormalTok{(}
\OperatorTok{:}\NormalTok{griddims }\OperatorTok{=\textgreater{}}\NormalTok{ (}\FloatTok{30}\NormalTok{, }\FloatTok{30}\NormalTok{),}
\OperatorTok{:}\NormalTok{max\_age }\OperatorTok{=\textgreater{}}\NormalTok{ [}\FloatTok{25}\NormalTok{, }\FloatTok{40}\NormalTok{],}
\OperatorTok{:}\NormalTok{init\_white }\OperatorTok{=\textgreater{}}\NormalTok{ [}\FloatTok{0.2}\NormalTok{, }\FloatTok{0.8}\NormalTok{],}
\OperatorTok{:}\NormalTok{init\_black }\OperatorTok{=\textgreater{}} \FloatTok{0.2}\NormalTok{,}
\OperatorTok{:}\NormalTok{albedo\_white }\OperatorTok{=\textgreater{}} \FloatTok{0.75}\NormalTok{,}
\OperatorTok{:}\NormalTok{albedo\_black }\OperatorTok{=\textgreater{}} \FloatTok{0.25}\NormalTok{,}
\OperatorTok{:}\NormalTok{surface\_albedo }\OperatorTok{=\textgreater{}} \FloatTok{0.4}\NormalTok{,}
\OperatorTok{:}\NormalTok{solar\_change }\OperatorTok{=\textgreater{}} \FloatTok{0.005}\NormalTok{,}
\OperatorTok{:}\NormalTok{solar\_luminosity }\OperatorTok{=\textgreater{}} \FloatTok{1.0}\NormalTok{,}
\OperatorTok{:}\NormalTok{scenario }\OperatorTok{=\textgreater{}} \OperatorTok{:}\NormalTok{ramp,}
\OperatorTok{:}\NormalTok{seed }\OperatorTok{=\textgreater{}} \FloatTok{165}\NormalTok{,}
\NormalTok{)}
\CommentTok{\# Создаём список всех комбинаций}
\NormalTok{params\_list }\OperatorTok{=} \FunctionTok{dict\_list}\NormalTok{(param\_dict)}
\ControlFlowTok{for}\NormalTok{ params }\KeywordTok{in}\NormalTok{ params\_list}
\NormalTok{model }\OperatorTok{=} \FunctionTok{daisyworld}\NormalTok{(;params}\OperatorTok{...}\NormalTok{)}
\FunctionTok{temperature}\NormalTok{(model) }\OperatorTok{=}\NormalTok{ StatsBase.}\FunctionTok{mean}\NormalTok{(model.temperature)}
\NormalTok{mdata }\OperatorTok{=}\NormalTok{ [temperature, }\OperatorTok{:}\NormalTok{solar\_luminosity]}
\NormalTok{agent\_df, model\_df }\OperatorTok{=} \FunctionTok{run!}\NormalTok{(model, }\FloatTok{1000}\NormalTok{; adata }\OperatorTok{=}\NormalTok{ adata, mdata }\OperatorTok{=}\NormalTok{ mdata)}
\NormalTok{figure }\OperatorTok{=}\NormalTok{ CairoMakie.}\FunctionTok{Figure}\NormalTok{(size }\OperatorTok{=}\NormalTok{ (}\FloatTok{600}\NormalTok{, }\FloatTok{600}\NormalTok{));}
\NormalTok{ax1 }\OperatorTok{=}\NormalTok{ figure[}\FloatTok{1}\NormalTok{, }\FloatTok{1}\NormalTok{] }\OperatorTok{=} \FunctionTok{Axis}\NormalTok{(figure, ylabel }\OperatorTok{=} \StringTok{"daisy count"}\NormalTok{)}
\NormalTok{blackl }\OperatorTok{=} \FunctionTok{lines!}\NormalTok{(ax1, agent\_df[!, }\OperatorTok{:}\NormalTok{time], agent\_df[!, }\OperatorTok{:}\NormalTok{count\_black], color }\OperatorTok{=} \OperatorTok{:}\NormalTok{red)}
\NormalTok{whitel }\OperatorTok{=} \FunctionTok{lines!}\NormalTok{(ax1, agent\_df[!, }\OperatorTok{:}\NormalTok{time], agent\_df[!, }\OperatorTok{:}\NormalTok{count\_white], color }\OperatorTok{=} \OperatorTok{:}\NormalTok{blue)}
\NormalTok{figure[}\FloatTok{1}\NormalTok{, }\FloatTok{2}\NormalTok{] }\OperatorTok{=} \FunctionTok{Legend}\NormalTok{(figure, [blackl, whitel], [}\StringTok{"black"}\NormalTok{, }\StringTok{"white"}\NormalTok{])}
\NormalTok{ax2 }\OperatorTok{=}\NormalTok{ figure[}\FloatTok{2}\NormalTok{, }\FloatTok{1}\NormalTok{] }\OperatorTok{=} \FunctionTok{Axis}\NormalTok{(figure, ylabel }\OperatorTok{=} \StringTok{"temperature"}\NormalTok{)}
\NormalTok{ax3 }\OperatorTok{=}\NormalTok{ figure[}\FloatTok{3}\NormalTok{, }\FloatTok{1}\NormalTok{] }\OperatorTok{=} \FunctionTok{Axis}\NormalTok{(figure, xlabel }\OperatorTok{=} \StringTok{"tick"}\NormalTok{, ylabel }\OperatorTok{=} \StringTok{"luminosity"}\NormalTok{)}
\FunctionTok{lines!}\NormalTok{(ax2, model\_df[!, }\OperatorTok{:}\NormalTok{time], model\_df[!, }\OperatorTok{:}\NormalTok{temperature], color }\OperatorTok{=} \OperatorTok{:}\NormalTok{red)}
\FunctionTok{lines!}\NormalTok{(ax3, model\_df[!, }\OperatorTok{:}\NormalTok{time], model\_df[!, }\OperatorTok{:}\NormalTok{solar\_luminosity], color }\OperatorTok{=} \OperatorTok{:}\NormalTok{red)}
\ControlFlowTok{for}\NormalTok{ ax }\KeywordTok{in}\NormalTok{ (ax1, ax2); ax.xticklabelsvisible }\OperatorTok{=} \ConstantTok{false}\NormalTok{; }\ControlFlowTok{end}
\NormalTok{plt\_name }\OperatorTok{=} \FunctionTok{savename}\NormalTok{(}\StringTok{"daisy{-}luminosity"}\NormalTok{,params) }\OperatorTok{*} \StringTok{".png"}
\FunctionTok{save}\NormalTok{(}\FunctionTok{plotsdir}\NormalTok{(plt\_name), figure)}
\ControlFlowTok{end}
\end{Highlighting}
\end{Shaded}

\begin{verbatim}
┌ Warning: Found `resolution` in the theme when creating a `Scene`. The `resolution` keyword for `Scene`s and `Figure`s has been deprecated. Use `Figure(; size = ...` or `Scene(; size = ...)` instead, which better reflects that this is a unitless size and not a pixel resolution. The key could also come from `set_theme!` calls or related theming functions.
└ @ Makie ~/.julia/packages/Makie/kJl0u/src/scenes.jl:264
┌ Warning: Found `resolution` in the theme when creating a `Scene`. The `resolution` keyword for `Scene`s and `Figure`s has been deprecated. Use `Figure(; size = ...` or `Scene(; size = ...)` instead, which better reflects that this is a unitless size and not a pixel resolution. The key could also come from `set_theme!` calls or related theming functions.
└ @ Makie ~/.julia/packages/Makie/kJl0u/src/scenes.jl:264
┌ Warning: Found `resolution` in the theme when creating a `Scene`. The `resolution` keyword for `Scene`s and `Figure`s has been deprecated. Use `Figure(; size = ...` or `Scene(; size = ...)` instead, which better reflects that this is a unitless size and not a pixel resolution. The key could also come from `set_theme!` calls or related theming functions.
└ @ Makie ~/.julia/packages/Makie/kJl0u/src/scenes.jl:264
┌ Warning: Found `resolution` in the theme when creating a `Scene`. The `resolution` keyword for `Scene`s and `Figure`s has been deprecated. Use `Figure(; size = ...` or `Scene(; size = ...)` instead, which better reflects that this is a unitless size and not a pixel resolution. The key could also come from `set_theme!` calls or related theming functions.
└ @ Makie ~/.julia/packages/Makie/kJl0u/src/scenes.jl:264
\end{verbatim}

\begin{Shaded}
\begin{Highlighting}[]
\ImportTok{using} \BuiltInTok{DrWatson}
\PreprocessorTok{@quickactivate} \StringTok{"project"}
\ImportTok{using} \BuiltInTok{Agents}
\ImportTok{using} \BuiltInTok{DataFrames}
\ImportTok{using} \BuiltInTok{Plots}
\ImportTok{using} \BuiltInTok{CairoMakie}
\FunctionTok{include}\NormalTok{(}\StringTok{"/home/srluipp/project/src/daisyworld.jl"}\NormalTok{)}
\CommentTok{\# Параметры эксперимента}
\NormalTok{param\_dict }\OperatorTok{=} \FunctionTok{Dict}\NormalTok{(}
\OperatorTok{:}\NormalTok{griddims }\OperatorTok{=\textgreater{}}\NormalTok{ (}\FloatTok{30}\NormalTok{, }\FloatTok{30}\NormalTok{),}
\OperatorTok{:}\NormalTok{max\_age }\OperatorTok{=\textgreater{}}\NormalTok{ [}\FloatTok{25}\NormalTok{, }\FloatTok{40}\NormalTok{],}
\OperatorTok{:}\NormalTok{init\_white }\OperatorTok{=\textgreater{}}\NormalTok{ [}\FloatTok{0.2}\NormalTok{, }\FloatTok{0.8}\NormalTok{],}
\OperatorTok{:}\NormalTok{init\_black }\OperatorTok{=\textgreater{}} \FloatTok{0.2}\NormalTok{,}
\OperatorTok{:}\NormalTok{albedo\_white }\OperatorTok{=\textgreater{}} \FloatTok{0.75}\NormalTok{,}
\OperatorTok{:}\NormalTok{albedo\_black }\OperatorTok{=\textgreater{}} \FloatTok{0.25}\NormalTok{,}
\OperatorTok{:}\NormalTok{surface\_albedo }\OperatorTok{=\textgreater{}} \FloatTok{0.4}\NormalTok{,}
\OperatorTok{:}\NormalTok{solar\_change }\OperatorTok{=\textgreater{}} \FloatTok{0.005}\NormalTok{,}
\OperatorTok{:}\NormalTok{solar\_luminosity }\OperatorTok{=\textgreater{}} \FloatTok{1.0}\NormalTok{,}
\OperatorTok{:}\NormalTok{scenario }\OperatorTok{=\textgreater{}} \OperatorTok{:}\NormalTok{default,}
\OperatorTok{:}\NormalTok{seed }\OperatorTok{=\textgreater{}} \FloatTok{165}\NormalTok{,}
\NormalTok{)}
\CommentTok{\# Создаём список всех комбинаций}
\NormalTok{params\_list }\OperatorTok{=} \FunctionTok{dict\_list}\NormalTok{(param\_dict)}
\ControlFlowTok{for}\NormalTok{ params }\KeywordTok{in}\NormalTok{ params\_list}
\NormalTok{model }\OperatorTok{=} \FunctionTok{daisyworld}\NormalTok{(;params}\OperatorTok{...}\NormalTok{)}
\FunctionTok{daisycolor}\NormalTok{(a}\OperatorTok{::}\DataTypeTok{Daisy}\NormalTok{) }\OperatorTok{=}\NormalTok{ a.breed}
\NormalTok{plotkwargs }\OperatorTok{=}\NormalTok{ (}
\NormalTok{agent\_color}\OperatorTok{=}\NormalTok{daisycolor,}
\NormalTok{agent\_size }\OperatorTok{=} \FloatTok{20}\NormalTok{,}
\NormalTok{agent\_marker }\OperatorTok{=} \CharTok{\textquotesingle{}✿\textquotesingle{}}\NormalTok{,}
\NormalTok{heatarray }\OperatorTok{=} \OperatorTok{:}\NormalTok{temperature,}
\NormalTok{heatkwargs }\OperatorTok{=}\NormalTok{ (colorrange }\OperatorTok{=}\NormalTok{ (}\OperatorTok{{-}}\FloatTok{20}\NormalTok{, }\FloatTok{60}\NormalTok{),),}
\NormalTok{)}
\NormalTok{plt1, \_ }\OperatorTok{=} \FunctionTok{abmplot}\NormalTok{(model; plotkwargs}\OperatorTok{...}\NormalTok{)}
\FunctionTok{step!}\NormalTok{(model, }\FloatTok{5}\NormalTok{)}
\NormalTok{plt2, \_ }\OperatorTok{=} \FunctionTok{abmplot}\NormalTok{(model; heatarray }\OperatorTok{=}\NormalTok{ model.temperature, plotkwargs}\OperatorTok{...}\NormalTok{)}
\FunctionTok{step!}\NormalTok{(model, }\FloatTok{40}\NormalTok{)}
\NormalTok{plt3, \_ }\OperatorTok{=} \FunctionTok{abmplot}\NormalTok{(model; heatarray }\OperatorTok{=}\NormalTok{ model.temperature, plotkwargs}\OperatorTok{...}\NormalTok{)}
\NormalTok{plt1\_name }\OperatorTok{=} \FunctionTok{savename}\NormalTok{(}\StringTok{"daisyworld"}\NormalTok{,params) }\OperatorTok{*} \StringTok{"\_step01"} \OperatorTok{*} \StringTok{".png"}
\NormalTok{plt2\_name }\OperatorTok{=} \FunctionTok{savename}\NormalTok{(}\StringTok{"daisyworld"}\NormalTok{,params) }\OperatorTok{*} \StringTok{"\_step04"} \OperatorTok{*} \StringTok{".png"}
\NormalTok{plt3\_name }\OperatorTok{=} \FunctionTok{savename}\NormalTok{(}\StringTok{"daisyworld"}\NormalTok{,params) }\OperatorTok{*} \StringTok{"\_step40"} \OperatorTok{*} \StringTok{".png"}
\FunctionTok{save}\NormalTok{(}\FunctionTok{plotsdir}\NormalTok{(plt1\_name), plt1)}
\FunctionTok{save}\NormalTok{(}\FunctionTok{plotsdir}\NormalTok{(plt2\_name), plt2)}
\FunctionTok{save}\NormalTok{(}\FunctionTok{plotsdir}\NormalTok{(plt3\_name), plt3)}
\ControlFlowTok{end}
\end{Highlighting}
\end{Shaded}

\begin{verbatim}
┌ Warning: Found `resolution` in the theme when creating a `Scene`. The `resolution` keyword for `Scene`s and `Figure`s has been deprecated. Use `Figure(; size = ...` or `Scene(; size = ...)` instead, which better reflects that this is a unitless size and not a pixel resolution. The key could also come from `set_theme!` calls or related theming functions.
└ @ Makie ~/.julia/packages/Makie/kJl0u/src/scenes.jl:264
┌ Warning: Found `resolution` in the theme when creating a `Scene`. The `resolution` keyword for `Scene`s and `Figure`s has been deprecated. Use `Figure(; size = ...` or `Scene(; size = ...)` instead, which better reflects that this is a unitless size and not a pixel resolution. The key could also come from `set_theme!` calls or related theming functions.
└ @ Makie ~/.julia/packages/Makie/kJl0u/src/scenes.jl:264
┌ Warning: Found `resolution` in the theme when creating a `Scene`. The `resolution` keyword for `Scene`s and `Figure`s has been deprecated. Use `Figure(; size = ...` or `Scene(; size = ...)` instead, which better reflects that this is a unitless size and not a pixel resolution. The key could also come from `set_theme!` calls or related theming functions.
└ @ Makie ~/.julia/packages/Makie/kJl0u/src/scenes.jl:264
┌ Warning: Found `resolution` in the theme when creating a `Scene`. The `resolution` keyword for `Scene`s and `Figure`s has been deprecated. Use `Figure(; size = ...` or `Scene(; size = ...)` instead, which better reflects that this is a unitless size and not a pixel resolution. The key could also come from `set_theme!` calls or related theming functions.
└ @ Makie ~/.julia/packages/Makie/kJl0u/src/scenes.jl:264
┌ Warning: Found `resolution` in the theme when creating a `Scene`. The `resolution` keyword for `Scene`s and `Figure`s has been deprecated. Use `Figure(; size = ...` or `Scene(; size = ...)` instead, which better reflects that this is a unitless size and not a pixel resolution. The key could also come from `set_theme!` calls or related theming functions.
└ @ Makie ~/.julia/packages/Makie/kJl0u/src/scenes.jl:264
┌ Warning: Found `resolution` in the theme when creating a `Scene`. The `resolution` keyword for `Scene`s and `Figure`s has been deprecated. Use `Figure(; size = ...` or `Scene(; size = ...)` instead, which better reflects that this is a unitless size and not a pixel resolution. The key could also come from `set_theme!` calls or related theming functions.
└ @ Makie ~/.julia/packages/Makie/kJl0u/src/scenes.jl:264
┌ Warning: Found `resolution` in the theme when creating a `Scene`. The `resolution` keyword for `Scene`s and `Figure`s has been deprecated. Use `Figure(; size = ...` or `Scene(; size = ...)` instead, which better reflects that this is a unitless size and not a pixel resolution. The key could also come from `set_theme!` calls or related theming functions.
└ @ Makie ~/.julia/packages/Makie/kJl0u/src/scenes.jl:264
┌ Warning: Found `resolution` in the theme when creating a `Scene`. The `resolution` keyword for `Scene`s and `Figure`s has been deprecated. Use `Figure(; size = ...` or `Scene(; size = ...)` instead, which better reflects that this is a unitless size and not a pixel resolution. The key could also come from `set_theme!` calls or related theming functions.
└ @ Makie ~/.julia/packages/Makie/kJl0u/src/scenes.jl:264
┌ Warning: Found `resolution` in the theme when creating a `Scene`. The `resolution` keyword for `Scene`s and `Figure`s has been deprecated. Use `Figure(; size = ...` or `Scene(; size = ...)` instead, which better reflects that this is a unitless size and not a pixel resolution. The key could also come from `set_theme!` calls or related theming functions.
└ @ Makie ~/.julia/packages/Makie/kJl0u/src/scenes.jl:264
┌ Warning: Found `resolution` in the theme when creating a `Scene`. The `resolution` keyword for `Scene`s and `Figure`s has been deprecated. Use `Figure(; size = ...` or `Scene(; size = ...)` instead, which better reflects that this is a unitless size and not a pixel resolution. The key could also come from `set_theme!` calls or related theming functions.
└ @ Makie ~/.julia/packages/Makie/kJl0u/src/scenes.jl:264
┌ Warning: Found `resolution` in the theme when creating a `Scene`. The `resolution` keyword for `Scene`s and `Figure`s has been deprecated. Use `Figure(; size = ...` or `Scene(; size = ...)` instead, which better reflects that this is a unitless size and not a pixel resolution. The key could also come from `set_theme!` calls or related theming functions.
└ @ Makie ~/.julia/packages/Makie/kJl0u/src/scenes.jl:264
┌ Warning: Found `resolution` in the theme when creating a `Scene`. The `resolution` keyword for `Scene`s and `Figure`s has been deprecated. Use `Figure(; size = ...` or `Scene(; size = ...)` instead, which better reflects that this is a unitless size and not a pixel resolution. The key could also come from `set_theme!` calls or related theming functions.
└ @ Makie ~/.julia/packages/Makie/kJl0u/src/scenes.jl:264
\end{verbatim}

\chapter{Выводы}\label{ux432ux44bux432ux43eux434ux44b}

При выполнении лабораторной работы я познакомилась с принципами
агентного моделирования, изучила модель daisyworld и построила графики
daisyworld в соответствии с принципами агентного моделирования.

\chapter*{Список
литературы}\label{ux441ux43fux438ux441ux43eux43a-ux43bux438ux442ux435ux440ux430ux442ux443ux440ux44b}
\addcontentsline{toc}{chapter}{Список литературы}

\printbibliography[heading=none]





\end{document}
